\section{Results}
\label{sec:experiments}

This section compares task formulations and training sources. Recall the dataset names: \textbf{LightwheelAI Baseline} (60 public full-task demonstrations), \textbf{Teleop} (manually collected subtasks), \textbf{Auto} (balanced auto-generated successful subtasks), and \textbf{Teleop+Auto} (their merge).

Every completed run uses the same 100-scene reference seed set (\texttt{pick\_orange\_reference\_100\_v1}) in the same order; the randomised quantities are the top-camera pose and the orange positions. Differences between models are therefore attributable to policy and formulation rather than to easier or harder sampled scenes.

\subsection{Full-Task Outcomes Across Formulations}

\Cref{fig:orange_outcome_distribution} reports final plate occupancy for every model on the shared 100-seed set under a 5000-step budget, via two metrics: the \emph{full-task success rate} (fraction of episodes ending with all three oranges placed --- whether the task is solved) and the \emph{mean final occupancy} (average oranges placed out of three --- crediting partial progress, separating headway from early stalls). Bars are grouped by training source and, within the two custom sources, by formulation (\emph{monotask} end-to-end vs.\ orchestrated \emph{subtask}).

\begin{figure}[H]
    \centering
    \includegraphics[width=\textwidth]{orange_outcome_distribution.pdf}
    \caption{Final plate occupancy across evaluated task formulations. Completed bars use the same 100 seeded initial conditions. Bar segments are percentages of evaluation episodes ending with 0, 1, 2, or 3 oranges in the plate; \textit{mean} is average final oranges in plate out of three. Bars are grouped by training dataset and labelled by task formulation.}
    \label{fig:orange_outcome_distribution}
\end{figure}

\paragraph{Policy backbone.} \textbf{[Placeholder: interpretation of the ACT vs.\ SmolVLA baseline comparison to be finalised.]}

\paragraph{Decomposition helps partial progress, not full-task completion.} On the same trajectories, the subtask formulation's benefit is in partial progress, not finishing all three. On Teleop it raises mean occupancy sharply (1.11 to 1.77) while leaving full-task success flat (18\% to 20\%): checkpoints convert empty and one-orange episodes into two-orange ones, rarely into three. On Teleop+Auto the monotask model is in fact stronger on both metrics (mean 1.50 vs.\ 1.42, full-task 26\% vs.\ 13\%). Decomposition thus secures more oranges on average on Teleop but does not raise full-task completion.

\paragraph{Subtasking rarely fails completely.} The clearest advantage is on the weakest outcome --- an episode with no orange placed. The subtask models almost always place at least one (97\% of Teleop episodes, 84\% on Teleop+Auto), whereas the monotask models leave many empty (only 57\% and 73\% reach one). Target redirection is the likely reason: a monotask policy can fixate on one hard-to-grasp orange and burn the episode, while the controller detects the repeated GRASP timeout and redirects to another. The decomposed formulation thus rarely fails totally, even when it does not complete the task.

\paragraph{Aggregating autonomous data did not help.} Adding the Auto episodes to Teleop regressed the subtask model on both metrics (full-task 20\% to 13\%, mean 1.77 to 1.42). In-distribution successful subtasks were thus not sufficient to improve full-task behaviour; the subtask-level analysis below isolates what degraded.

\subsection{Lift Failures and Recovery}
\label{sec:lift_recovery}

Beyond the final outcome, the per-step traces show how each policy fails when a held orange slips. A \textbf{lift} begins once a held orange rises above its start height; it is a \textbf{drop} if the orange leaves the gripper before the plate. \Cref{tab:lift_recovery} reports \textbf{lift success} --- successes / genuine lift attempts (successes + drops, excluding episode-end truncations). Whether a dropped orange is itself reclaimed (same-orange, Level-1 recovery) is deferred to the formulation-fair task-level view of \cref{sec:task_recovery}, since the monotask has no fixed target to return to and the comparison would otherwise favour the orchestrator by construction.

One caveat frames every subtask-level comparison below: the monotask traces are reconstructed offline from physics (\cref{subsec:failure_taxonomy}), whereas the subtask traces are the orchestrator's own log.

\begin{table}[H]
    \centering
    \small
    \caption{Lift success by training set and task formulation: successes / genuine lift attempts (successes + drops, excluding episode-end truncations).}
    \label{tab:lift_recovery}
    \begin{tabular}{llc}
        \toprule
        Training set & Formulation & Lift success \\
        \midrule
        Teleop      & Monotask & 60.7\% (142/234) \\
        Teleop      & Subtask  & 92.5\% (209/226) \\
        \midrule
        Teleop+Auto & Monotask & 83.2\% (183/220) \\
        Teleop+Auto & Subtask  & 88.8\% (174/196) \\
        \bottomrule
    \end{tabular}
\end{table}

\paragraph{The subtask policies lift more reliably.} Lift success is markedly higher under the subtask formulation (89--92\% vs.\ 61--83\% for monotask), most starkly on Teleop where the monotask model loses nearly two of every five held oranges. Part of this gap is by construction: the orchestrator enters LIFT only once its grasp check passes, so a subtask lift starts from a \emph{verified} solid grasp, whereas the monotask trace infers a lift from any upward motion of a held orange and counts lifts that begin from marginal grasps. The higher subtask rate is thus partly a better-conditioned starting point, not only a better policy.

\paragraph{When the monotask lift fails, it fails immediately.} The monotask drops are not mid-lift losses but \emph{immediate slips}: the held orange barely leaves the table before it is lost (\cref{tab:failure_geometry}), the signature of a grasp that was never solid. One subtask-specific wrinkle appears on Teleop+Auto, whose lift attempts also \emph{time out} 9.7\% of the time (21/217) without dropping the orange --- the policy, trained to freeze at each subtask boundary, sometimes idles on a securely held orange instead of completing the lift, an artefact of the tail-freeze training discussed in \cref{sec:discussion}.

\subsection{Placing Failures}
\label{sec:place}

Placing the lifted orange is the final subtask. A \textbf{place} succeeds when the orange comes to rest in the plate, read from the end-of-attempt geometry --- independent of gripper pose, and independent of whether the orchestrator confirmed the release within its step budget (the timeout caveat of \cref{sec:discussion}). Being purely geometric, the test applies uniformly to both formulations. \Cref{tab:place_success} reports the \textbf{place success rate} overall and for the \textbf{first} placement into an empty plate --- the one scene-state split that carries signal.

\begin{table}[H]
    \centering
    \small
    \caption{Place success rate by training set and formulation: the fraction of place attempts whose target orange is in the plate at the attempt's end (a per-attempt geometric read; full-episode truncations excluded). \emph{First placement} is the subset placing into an empty plate, where the monotask gap is concentrated.}
    \label{tab:place_success}
    \begin{tabular}{llcc}
        \toprule
        Training set & Formulation & Place success & First placement (empty plate) \\
        \midrule
        Teleop      & Monotask & 75.4\% (107/142) & 58.0\% (40/69) \\
        Teleop      & Subtask  & 89.0\% (186/209) & 88.4\% (99/112) \\
        \midrule
        Teleop+Auto & Monotask & 82.5\% (151/183) & 79.8\% (71/89) \\
        Teleop+Auto & Subtask  & 81.6\% (142/174) & 83.0\% (83/100) \\
        \bottomrule
    \end{tabular}
\end{table}

\paragraph{The subtask policies place reliably; the Teleop monotask struggles to start.} Once an orange is lifted, the subtask policies place it about four times in five (89\% Teleop, 82\% Teleop+Auto). The \textbf{Teleop monotask succeeds on only 58\% of \emph{first} placements} into an empty plate --- the difficulty is getting the first orange in, not a late collapse, mirroring its many empty episodes (\Cref{fig:orange_outcome_distribution}). \textbf{Placing is not the source of the subtask advantage}: once lifting succeeds both formulations finish comparably; they diverge earlier, at grasping and lifting an orange at all.

\paragraph{Why the monotask attempts fail.} For the monotask models, whose every attempt logs the gripper and orange geometry at the decisive moment, the two failure modes have a clean physical signature (\Cref{tab:failure_geometry}). A failed place releases the orange a median 0.16--0.21\,m from the plate centre --- four to five times the 0.04\,m of a success --- so the orange is dropped over the \emph{rim}, not into the plate. A failed lift gains only 3--4\,mm of height before the orange is lost, against about 70\,mm for a success: the grasp was never solid enough to raise it. These signatures appear \emph{only} under the monotask formulation, and that is itself informative. The subtask orchestrator enters LIFT only after a force-confirmed grasp, so the policy is never asked to lift or place from the marginal grasps that produce these immediate slips and rim releases: the guard removes exactly the regime in which the unguarded VLA, judging the moment itself, fails. (A failed place is also seldom recovered within the episode; the formulation-fair recovery view is Level-2, \cref{sec:task_recovery}.)

\begin{table}[H]
    \centering
    \small
    \caption{Failure geometry of the \textbf{monotask} policies (median over attempts; the subtask checkpoints do not log per-attempt geometry). \emph{Place release dist.}\ is the orange's horizontal distance from the plate centre when released --- a slip lets go over the rim. \emph{Lift height gain} is how far a held orange rose before it was lost --- a drop never really leaves the table. These unguarded modes surface only in the monotask: the subtask orchestrator gates LIFT on a force-confirmed grasp.}
    \label{tab:failure_geometry}
    \begin{tabular}{lcccc}
        \toprule
        & \multicolumn{2}{c}{Place release dist.\ (m)} & \multicolumn{2}{c}{Lift height gain (m)} \\
        \cmidrule(lr){2-3}\cmidrule(lr){4-5}
        Training set & success & slip & success & drop \\
        \midrule
        Teleop      & 0.04 & 0.16 & 0.067 & 0.004 \\
        Teleop+Auto & 0.04 & 0.21 & 0.074 & 0.003 \\
        \bottomrule
    \end{tabular}
\end{table}

\subsection{Task-Level Recovery}
\label{sec:task_recovery}

The lift and place analyses measure \textbf{Level-1 (same-orange) recovery}, the orchestrator's local retry --- which the monotask, never given a fixed target, has no mechanism to aim for. For an equal-footing comparison we also measure \textbf{Level-2 (task) recovery} (\cref{subsec:orchestrator}): after a setback, does the system keep scoring by placing \emph{some} orange, regardless of which? Because the objective is indifferent to which orange is placed, this is where monotask and subtask are genuinely comparable, needing no inference of intended target.

The trigger is the same setback for both formulations --- a LIFT drop or PLACE slip from simulator state --- and recovery is a later successful placement. \Cref{tab:task_recovery} reports two views. \emph{Score recovery}: of episodes with at least one setback, the fraction whose placed count rose afterward (same or different orange). \emph{Different-orange recovery}, the stricter Level-2 signal: of setbacks with at least one \emph{other} orange still unplaced, the fraction followed by placing a \emph{different} orange. The 5000-step budget is shared, so exhaustion is symmetric across formulations.

\begin{table}[H]
    \centering
    \small
    \caption{Level-2 (task) recovery after a physical setback (LIFT drop or PLACE slip), measured identically for both formulations from simulator state. \emph{Score recovery}: of episodes with at least one setback, the fraction whose placed count rose after the first setback. \emph{Different-orange recovery}: of setbacks with at least one other orange still unplaced, the fraction followed by a successful placement of a different orange. Denominators differ because the subtask policy suffers far fewer setbacks.}
    \label{tab:task_recovery}
    \begin{tabular}{llcc}
        \toprule
        Training set & Formulation & Score recovery & Different-orange recovery \\
        \midrule
        Teleop      & Monotask & 35.5\% (27/76) & 23.1\% (28/121) \\
        Teleop      & Subtask  & 73.5\% (25/34) & 60.0\% (18/30)  \\
        \midrule
        Teleop+Auto & Monotask & 53.7\% (22/41) & 35.1\% (20/57)  \\
        Teleop+Auto & Subtask  & 42.9\% (18/42) & 28.0\% (14/50)  \\
        \bottomrule
    \end{tabular}
\end{table}

\paragraph{Level-2 is where the orchestrator's advantage is real --- and where Auto data erodes it.} On Teleop the subtask formulation suffers far fewer setbacks (34 vs.\ 76 episodes) and recovers from them far more often: it keeps scoring in 73.5\% of episodes against 35.5\% for monotask, and reaches a \emph{different} orange after 60.0\% of setbacks against 23.1\%. This is the mechanism behind ``rarely fails completely'' (\cref{fig:orange_outcome_distribution}): Level-2 redirection turns dead-ends into progress the end-to-end policy, prone to fixating, does not make. The pattern reverses on Teleop+Auto, where the subtask model recovers \emph{less} often on both views (42.9\% vs.\ 53.7\%; 28.0\% vs.\ 35.1\%) --- the full-task regression now visible at the recovery level: the added data weakened exactly the task-level robustness the orchestrator exists to provide.

\subsection{Target Obedience in the Subtask Policies}
\label{sec:obedience}

Only the orchestrated runs name their target by spatial label (``Grasp \emph{<position>} orange''), so we can ask whether the policy grasps the orange it was \emph{told} to --- a question monotask runs cannot pose. A grasp is \textbf{obeyed} when the gripper closes on the requested orange (matched by identity, not position label), counted over grasps that engage an orange; any other engaged grasp is a \textbf{misgrab}. We report obedience by \textbf{scene state} (oranges already placed at grasp time), since resolving a referent depends on how many distractors remain.

\paragraph{The instructions are dominated by one referent.} \Cref{tab:label_prevalence} counts training GRASP episodes per label: \emph{right} alone is almost half, with \emph{middle} and \emph{left} next and diagonal labels rare; the Auto episodes only deepened this skew. This sets up a clean test: \textbf{if obedience were limited by exposure, the most common labels should be obeyed best.}

\begin{table}[H]
    \centering
    \small
    \caption{Prevalence of positional grasp instructions in the training data: number of GRASP episodes per label, by source. ``Other'' aggregates the rarely used top left, top middle, and bottom middle labels.}
    \label{tab:label_prevalence}
    \begin{tabular}{lccc}
        \toprule
        Position & Teleop & Auto & Combined \\
        \midrule
        top right    & 25  & 4   & 29  \\
        left         & 34  & 11  & 45  \\
        middle       & 74  & 48  & 122 \\
        right        & 135 & 74  & 209 \\
        bottom right & 15  & 1   & 16  \\
        other        & 7   & 5   & 12  \\
        \midrule
        Total        & 290 & 143 & 433 \\
        \bottomrule
    \end{tabular}
\end{table}

\paragraph{Obedience improves as the scene simplifies.} \Cref{fig:grasp_confusion} reports, per model and scene state, how often each requested position was obeyed (green diagonal) and where misses landed (red). Roughly three in four engaged grasps reach the requested orange, and the rate climbs as oranges clear: lowest with all three present, higher once one is placed, trivially perfect with two placed. \textbf{Fewer distractors make the referent easier to resolve.}

\begin{figure}[H]
    \centering
    \includegraphics[width=\textwidth]{grasp_obedience_confusion.pdf}
    \caption{Grasp-obedience confusion matrices for the two subtask models, with all three oranges on the table (top) and one already placed (bottom). Rows are the requested position, columns the position actually grasped (by orange identity). \textbf{Green} diagonal cells are obeyed grasps and \textbf{red} off-diagonal cells are misgrabs, each shaded by its share of the row; the per-row obeyed rate and count, and the per-panel overall rate, are printed alongside. Labels reflect the scene at grasp time. The 2-placed state (one orange left, trivially obeyed) and two single-sample requests are omitted.}
    \label{fig:grasp_confusion}
\end{figure}

\paragraph{Misgrabs follow a systematic inward bias.} The errors are far from random: in every panel of \Cref{fig:grasp_confusion} a missed request closes on the orange immediately \textbf{inward of} the intended one --- \emph{right} taken as \emph{middle}, \emph{middle} as \emph{left}, \emph{bottom right} as \emph{top right}. The extreme referents \emph{left} and \emph{top right} are thus obeyed most reliably, while \emph{right} is the worst-obeyed label in both models and scene states. This is real behaviour, not a logging artefact: the grasped orange is identified from contact, the requested one is verifiably the true rightmost, and the orange taken is a physically distinct one a few centimetres inward --- a genuine short reach in packed scenes.

\paragraph{A grounding gap, not a coverage gap.} This ordering runs \emph{opposite} to training prevalence: the most frequent label \emph{right} is the worst obeyed, the sparse \emph{top right} among the best. \textbf{More examples do not buy better obedience} --- exactly what the Auto data shows, adding exposure without lifting it. The bottleneck is \emph{grounding} the instruction to the scene: the VLA's positional conditioning is a spatially biased channel, not a clean target selector --- the context-dependence the scripted spatial reset works around.
